\section*{Hoofdstuk \ref{ch:b2:population-genetics} --- Populatiegenetica}

\begin{solution}{exo:b2:population-genetics:1}
$p(M) = (1280 + 320)/2000 = 0.8$, $q(N) = 0.2$. Verwacht $MM$:
$0.64\times 1000 = 640$; $MN$: $2\times 0.8\times 0.2\times 1000 =
320$; $NN$: 40 --- precies de waargenomen aantallen. De populatie is
voor dit locus in evenwicht.
\end{solution}

\begin{solution}{exo:b2:population-genetics:2}
$q = \sqrt{1/20000} = 0.0071$; dragers $2pq \approx 0.014$, één
persoon op 70. Dragers overtreffen getroffenen met $2p/q \approx
280$ tegen één.
\end{solution}

\begin{solution}{exo:b2:population-genetics:3}
\omterm{def:b2:population-genetics:fitness}{Fitness}: het verwachte aantal nakomelingen van een \omterm{def:b2:meiosis-heredity:vocabulary}{genotype} ten opzichte van
het beste. \omterm{def:b2:population-genetics:fitness}{Selectiecoëfficiënt} $s$: het fitnesstekort, $w = 1 -
s$. \omterm{thm:b2:population-genetics:heritability}{Erfelijkheidsgraad} $h^{2}$: de fractie van de fenotypische variantie
die additief-genetisch is, of gelijkwaardig de helling van nakomelingen op
ouders. Krachten: \omterm{def:b2:genome-diversification:mutation}{mutatie}, selectie, drift, migratie, niet-willekeurige
paring (de laatste verandert genotypefrequenties, geen \omterm{def:b2:population-genetics:frequencies}{allelfrequenties}).
\end{solution}

\begin{solution}{exo:b2:population-genetics:4}
Gericht: melanisme bij de berkenspanner, met als factor de vogels die op zicht
jagen op door roet verdonkerde stammen. Stabiliserend: het geboortegewicht
bij de mens, met als factor de sterfte van zeer kleine en zeer grote baby's.
Balancerend: sikkelcel, met als factoren malaria (tegen $AA$) en bloedarmoede
(tegen $SS$) samen.
\end{solution}

\begin{solution}{exo:b2:population-genetics:5}
$q_t = q_0/(1 + tq_0)$: halveren duurt $1/q_0 = 50$ generaties;
$0.001$ bereiken duurt $1/0.001 - 1/0.02 = 950$ generaties. Bijna
alle kopieën van het \omterm{def:b2:meiosis-heredity:vocabulary}{allel} zitten in heterozygoten, onzichtbaar voor
selectie; getroffenen verhinderen zich voort te planten (wat de natuur
al doet, want ze zijn dood) verandert binnen een mensenleven niets
meetbaars.
\end{solution}

\begin{solution}{exo:b2:population-genetics:6}
$p = 0.3$: $\bar w = 0.09 + 0.42 + 0.8\times 0.49 = 0.902$; $p' =
(0.09 + 0.21)/0.902 = 0.333$; $\Delta p = +0.033$. $p = 0.9$: $\bar w
= 0.998$, $p' = 0.9018$, $\Delta p = +0.0018$. Bij $p = 0.9$ is maar
$q^{2} = 0.01$ van de populatie $aa$: het \omterm{def:b2:meiosis-heredity:vocabulary}{allel} waartegen selectie werkt,
zit verborgen in heterozygoten, en $\Delta p = spq^{2}/\bar w$
bevat de factor $q^{2}$.
\end{solution}

\begin{solution}{exo:b2:population-genetics:7}
Met $s = 0.12$ tegen $AA$ en $t = 0.8$ tegen $SS$ is $\hat q =
s/(s + t) = 0.12/0.92 = 0.13$. Getroffen geboorten $\hat q^{2} = 0.017$,
ongeveer één kind op 60 --- de prijs van de bescherming die de
heterozygoten genieten.
\end{solution}

\begin{solution}{exo:b2:population-genetics:8}
\omterm{def:b2:meiosis-heredity:vocabulary}{Recessief} letaal: $\hat q = \sqrt{\mu} = \sqrt{2\times 10^{-5}} =
0.0045$; getroffen geboorten $\hat q^{2} = \mu = 2\times 10^{-5}$, één op
\num{50000}. \omterm{def:b2:meiosis-heredity:vocabulary}{Dominant} letaal: elke kopie wordt verwijderd in de
generatie waarin ze verschijnt, zodat de getroffen geboorten gewoon de nieuwe
\omterm{def:b2:genome-diversification:mutation}{mutaties} zijn, $2\mu = 4\times 10^{-5}$ (twee gameten per geboorte), één op
\num{25000} --- het recessieve \omterm{def:b2:meiosis-heredity:vocabulary}{allel} verbergt zijn kopieën tweehonderd tegen
één, het dominante verbergt niets.
\end{solution}

\begin{solution}{exo:b2:population-genetics:9}
$H_t = 0.5\,(1 - 1/100)^{t}$: na 10 generaties $0.45$; na 50
$0.30$; na 200 $0.067$. Voor $(1 - 1/2N)^{100} = 0.9$:
$100/(2N) \approx -\ln 0.9 = 0.105$, dus $N \approx 475$, ongeveer 500
zich voortplantende individuen.
\end{solution}

\begin{solution}{exo:b2:population-genetics:10}
Twintig genkopieën: $P(\text{afwezig}) = 0.9^{20} = 0.12$. Eén enkele
kopie onder 20 heeft frequentie $0.05$, en een neutraal \omterm{def:b2:meiosis-heredity:vocabulary}{allel} raakt gefixeerd
met een kans gelijk aan zijn frequentie: $0.05$. Op een vasteland van, zeg,
\num{100000} vogels heeft één kopie een kans van $5\times 10^{-6}$
om gefixeerd te raken. Het eiland is een plek waar zeldzame \omterm{def:b2:meiosis-heredity:vocabulary}{allelen} verloren
gaan en waar de \omterm{def:b2:meiosis-heredity:vocabulary}{allelen} die overleven de overhand kunnen krijgen: drift
verarmt en differentieert tegelijk.
\end{solution}

\begin{solution}{exo:b2:population-genetics:11}
$R = h^{2}S = 0.3\times 1.4\times \qty{1000}{L} = \qty{420}{L}$ per
generatie. Met stieren uit de beste \qty{1}{\%} is het gemiddelde
differentieel $(1.4 + 2.7)/2 = 2.05$ standaardafwijkingen en $R =
0.3\times 2050 = \qty{615}{L}$. De respons vertraagt omdat selectie
de additieve variantie opgebruikt --- gunstige \omterm{def:b2:meiosis-heredity:vocabulary}{allelen} raken gefixeerd
en $h^{2}$ daalt --- en omdat de gecorreleerde veranderingen in andere
eigenschappen (vruchtbaarheid, gezondheid) kosten opleggen die selectie op
alleen de melkgift niet ziet.
\end{solution}

\begin{solution}{exo:b2:population-genetics:12}
Selectie ziet \omterm{def:b2:meiosis-heredity:vocabulary}{fenotypen}. Een \omterm{def:b2:meiosis-heredity:vocabulary}{recessief} \omterm{def:b2:meiosis-heredity:vocabulary}{allel} in een \omterm{def:b2:meiosis-heredity:vocabulary}{heterozygoot} geeft geen
\omterm{def:b2:meiosis-heredity:vocabulary}{fenotype}, dus selectie kan het niet verwijderen, en zijn frequentie daalt als
$1/t$, steeds trager. Een neutrale \omterm{def:b2:genome-diversification:mutation}{mutatie} maakt geen verschil in \omterm{def:b2:population-genetics:fitness}{fitness};
haar lot hangt alleen van drift af, en ze raakt gefixeerd met kans $1/2N$. Een
kenmerk met een \omterm{thm:b2:population-genetics:heritability}{erfelijkheidsgraad} nul
kan sterk worden geselecteerd --- de overlevenden kunnen sterk van de
populatie verschillen --- en toch is de volgende generatie onveranderd, omdat
de verschillen niet erfelijk waren. Selectie kan alleen inwerken op erfelijke
verschillen die voor de \omterm{def:b2:population-genetics:fitness}{fitness} verschil maken; al het andere evolueert door
toeval of helemaal niet.
\end{solution}

\begin{solution}{pb:b2:population-genetics:1}
\textbf{1.} Melanistische fractie $1 - q^{2} = 1 - 0.999^{2} = 0.002$,
één vlinder op 500; de fractie van de $M$-allelen in heterozygoten is
$2pq/(2pq + 2p^{2}) = q = 0.999$: vrijwel alle.
\textbf{2.} Met $\bar w = 1 - sq^{2}$ zijn de kopieën van het lichte \omterm{def:b2:meiosis-heredity:vocabulary}{allel}
na selectie $pq\cdot 1 + q^{2}(1 - s)$ op $\bar w$, dus $q'
= (pq + q^{2} - sq^{2})/(1 - sq^{2}) = q(1 - sq)/(1 - sq^{2})$.
\textbf{3.} $q_1 = 0.999\times(1 - 0.2997)/(1 - 0.2994) = 0.99857$:
een verandering van $-0.0014$.
\textbf{4.} $q = \sqrt{0.1} = 0.32$.
\textbf{5.} De herhaalde recursie neemt 28 generaties voor $s =
0.3$; de benadering geeft dezelfde orde (de verandering is eerst
$-0.0003$ en versnelt daarna naarmate $p$ groeit). De waargenomen 50
generaties komen overeen met een kleinere $s$, ongeveer $0.2$ (de recursie
geeft 45 generaties voor $s = 0.2$): een selectie van \qtyrange{20}{30}{\%}
per generatie, enorm naar de maatstaven van de meeste \omterm{def:b2:meiosis-heredity:vocabulary}{allelen}.
\textbf{6.} Met $w_{MM} = w_{Mm} = 1 - s$, $w_{mm} = 1$: $\bar w = 1
- s(1 - q^{2})$ en $p' = p(1 - s)/\bar w$, dus $\Delta p = p[(1 - s) -
\bar w]/\bar w = -spq^{2}/\bar w$. De lichte vorm wordt alleen als
\omterm{def:b2:meiosis-heredity:vocabulary}{homozygoot} $mm$ begunstigd, aan het begin een fractie $q^{2} = 0.1$: de
terugkeer begint traag. Als $q \to 1$ is $\Delta p \approx -sp$ en neemt $M$
exponentieel af, met een factor $1 - s$ per generatie: een
\omterm{def:b2:meiosis-heredity:vocabulary}{dominant} \omterm{def:b2:meiosis-heredity:vocabulary}{allel} kan zich nergens verbergen. De opkomst was het spiegelbeeld:
$\Delta p \approx +sp$ zolang $M$ zeldzaam was (exponentieel begin), daarna
$\Delta q \approx -sq^{2}p \to$ een slakkengang, doordat het lichte \omterm{def:b2:meiosis-heredity:vocabulary}{allel} zich
in heterozygoten verborg. De recursie geeft 27 generaties voor de terugkeer
van $p = 0.68$ tot $0.001$, waarvan 16 om een frequentie van \qty{5}{\%} van
het melanistische \omterm{def:b2:meiosis-heredity:vocabulary}{allel} te bereiken.
\textbf{7.} $24$ genkopieën: $0.95^{24} = 0.29$.
\textbf{8.} $H_1 = 0.4\,(1 - 1/24) = 0.383$.
\textbf{9.} $H_{100} = 0.383\,(1 - 1/80)^{100} = 0.383\times 0.284 =
0.109$: \qty{27}{\%} van de $0.4$ van het vasteland.
\textbf{10.} $P(\text{fixatie}) = 1/(2N) = 1/80 = 0.0125$; verwachte
tijd tot fixatie $4N = 160$ generaties.
\textbf{11.} Zonder inteelt $q^{2} = 0.0025$; met $F = 0.3$
$q^{2} + Fpq = 0.0025 + 0.3\times 0.95\times 0.05 = 0.017$: zeven
keer meer getroffen vogels.
\textbf{12.} Eén migrant per generatie ($m = 1/40$) volstaat om te
voorkomen dat het eiland op neutrale loci uiteenloopt, en om de \omterm{def:b2:meiosis-heredity:vocabulary}{allelen}
aan te vullen die door drift verloren zijn gegaan; hij brengt ook de
schadelijke \omterm{def:b2:meiosis-heredity:vocabulary}{allelen} van het vasteland binnen met hun frequentie op het
vasteland, en die worden door de inteelt op het eiland als homozygoten
zichtbaar --- de last wordt bepaald door het evenwicht tussen immigratie en
blootstelling.
\textbf{13.} $\hat q = \sqrt{10^{-5}} = 0.0032$; getroffen geboorten $\hat
q^{2} = 10^{-5}$.
\textbf{14.} $\hat q = \sqrt{10^{-5}/0.1} = 0.01$; getroffen geboorten
$10^{-4}$: een tien keer zwakkere selectie laat het \omterm{def:b2:meiosis-heredity:vocabulary}{allel} drie keer zo
algemeen worden en de ziekte tien keer.
\textbf{15.} $\hat p = \mu/(hs) = 10^{-5}/0.05 = 2\times 10^{-4}$;
dragers $2\hat p = 4\times 10^{-4}$.
\textbf{16.} Van $\hat q = 0.0032$ naar $0.01$: het \omterm{def:b2:meiosis-heredity:vocabulary}{allel} stijgt met een
snelheid die door \omterm{def:b2:genome-diversification:mutation}{mutatie} wordt bepaald, $\mu = 10^{-5}$ per generatie, dus de
nadering duurt in de orde van $\Delta q/\mu \approx 700$ generaties ---
\num{20000} jaar: elke beslissing die nu wordt genomen, wordt gevoeld door een
populatie die in niets op de onze zal lijken.
\textbf{17.} Elk locus levert $2\hat q = 0.0063$ letale kopieën
per persoon; over \num{20000} loci ongeveer $130$ --- maar dat veronderstelt
dat elk gen tot een \omterm{def:b2:meiosis-heredity:vocabulary}{recessief} \omterm{prop:b2:meiosis-heredity:extensions}{letaal allel} kan muteren; met de gangbare
schatting van enkele duizenden essentiële genen draagt een persoon enkele tot
enkele tientallen letale equivalenten in heterozygote toestand.
\textbf{18.} Elk zit op een ander locus en elk is zeldzaam, zodat de
kans dat een willekeurig paar allebei hetzelfde draagt klein is;
neven en nichten delen een achtste van hun genen door afstamming, zodat voor
elk van de verscheidene letale \omterm{def:b2:meiosis-heredity:vocabulary}{allelen} van een neef het risico dat het kind
\omterm{def:b2:meiosis-heredity:vocabulary}{homozygoot} is $1/16$ per gedeeld \omterm{prop:b2:meiosis-heredity:extensions}{letaal allel} bedraagt --- de extra
kindersterfte bij huwelijken tussen neef en nicht is gemeten op enkele procenten.
\textbf{19.} $S = 1.16\times \qty{800}{L} = \qty{930}{L}$; $R = 0.4\times
930 = \qty{370}{L}$.
\textbf{20.} Vijf generaties: ongeveer \qty{1860}{L}; in de praktijk minder,
omdat de additieve variantie uitgeput raakt, omdat $h^{2}$
in de oorspronkelijke kudde werd geschat, en omdat de omgeving
(voer, huisvesting) die de beste koeien voortbracht niet wordt doorgegeven.
\textbf{21.} Gemiddeld differentieel $(1.16 + 2.4)/2 = 1.78$
standaardafwijkingen, $S = \qty{1424}{L}$, $R = 0.4\times 1424 = \qty{570}{L}$.
\textbf{22.} Klonen hebben geen genetische variantie, dus $h^{2} = 0$ en $R =
0$, wat $S$ ook is: \omterm{thm:b2:population-genetics:heritability}{erfelijkheidsgraad} is een eigenschap van een populatie,
het deel van haar variantie dat genetisch is, geen eigenschap van het kenmerk.
\textbf{23.} $R = 0.7\times 0.6 = \qty{0.42}{mm}$ dieper.
\textbf{24.} Dezelfde respons $\qty{370}{L}$ vraagt $S = 370/0.1 =
\qty{3700}{L}$, dat is $4.6$ standaardafwijkingen --- minder dan één koe op
tienduizend houden: onpraktisch, en daarom worden eigenschappen met een lage
\omterm{thm:b2:population-genetics:heritability}{erfelijkheidsgraad} traag of met andere middelen verbeterd
(nakomelingenonderzoek, grotere families).
\textbf{25.} Vlinder: $s \approx 0.2$ tot $0.3$, 28 generaties bij $0.3$
en 45 bij $0.2$, tegenover 50 waargenomen; eiland: $H_{100} \approx 0.11$,
\qty{27}{\%} van het vasteland; kudde: \qty{370}{L} per generatie met
alleen de koeien, \qty{570}{L} met geselecteerde stieren.
\end{solution}
